\documentclass[../../main2.tex]{subfiles}

\begin{document}

	% ======================================================================
	% 骨架文件：小节标题与追问链为终版；正文由后续会话填充。
	% 扩增模式标注：本章 = 有压扩增 + 身份权重（部分守恒）。
	% 原书材料接入点：
	%   - 原书 part3 ch3（政治与权力结构）→ 本章主体
	%   - 原书 part2 ch2（生存策略分化）→ §11.4：集权 vs 碎片改写为
	%     两种有压扩增组织形态（从「地理结果」改为「扩增模式」）
	% ======================================================================

	\chapter{通道三：政治（经济 + 身份权重）}
	\label{ch:politics}

	\begin{motivation}
		上一章遗留的问题：经济交换在模型里是对称的等价交换——但交换者有身份、历史、偏好。当交换双方的身份关系变得重要，交换网络会变成什么？\\
		\textbf{核心动机}：论证政治从经济中分离的时刻——交换中加入不可货币化的偏好；权力的来源是对职能的集体认同；解释政治守恒性为何弱于经济。\\
		\textbf{扩增模式}：有压 + 部分守恒——资源基础守恒，认同可以凭空增减。
	\end{motivation}

	\begin{logicmap}
	\begin{tikzpicture}[node distance=0.9cm, every node/.style={draw, rectangle, rounded corners, align=center}]
	\node (q) {交换双方的身份变重要，会怎样？};
	\node (i) [below=of q] {从等价到不等价（§11.1）};
	\node (r) [below=of i] {认同赋予权力（§11.2）};
	\node (c) [below=of r] {权力的模因化 = 制度化（§11.3）};
	\node (s) [below=of c] {生存策略的两条扩增路径（§11.4）};
	\node (n) [below=of s] {权力需要约束（抛给第12章）};
	\draw[-{Stealth}] (q) -- (i);
	\draw[-{Stealth}] (i) -- (r);
	\draw[-{Stealth}] (r) -- (c);
	\draw[-{Stealth}] (c) -- (s);
	\draw[-{Stealth}] (s) -- (n);
	\end{tikzpicture}
	\end{logicmap}

	% ==================================================================
	\section{从等价到不等价}
	\label{sec:pol-nonequivalence}

	如果交换只认货币尺度，定价就是全部信息：一条烟值多少钱，标在货架上，谁买都一样。但现在加入一个约束：交换携带身份。同一条烟，你在商店买是两块钱，你岳父家抽的那条没有价格——它是一句「爸，给您带的」。人情、义务、面子，这些东西不在货币尺度里，但它们真实地改变着交换的条件。

	这一步改变了所有性质。当交换者不再是匿名原子而是有身份关系的个体时，「等价」这个概念就失效了：你和老板谈薪资是等价交换（你的劳动换他的钱），但你帮老板搬家不是——你得不到等价的回报，因为你们之间的关系不是市场关系。政治就是从这种「定价失效」中分离出来的：凡是无法用货币度量的交换偏好，都进入了政治通道。

	\begin{readingnote}
		给邻居送钱是侮辱，帮邻居搬家是情分。为什么？因为送钱把你们的关系拉进了经济通道（等价交换），而搬家是在政治通道里操作（身份关系中的义务表达）。两种通道的规则不同，所以同一个动作在不同通道里有完全不同的含义。
	\end{readingnote}

	\paragraph*{逻辑衔接}
	定价失效之后，新的度量出现了——别人为什么听你的？下一节回答这个问题。

	% ==================================================================
	\section{认同赋予权力}
	\label{sec:pol-recognition}
	权力的来源不是暴力，是对职能的集体认同。一个将军之所以能指挥军队，不是因为他的枪比士兵多，而是因为整个组织认同「将军有权发号施令」这个事实。认同是权力的燃料——没有它，命令就是空话。

	这解释了为什么政治守恒性弱于经济。经济通道里，价值总量守恒：你赚一块钱，有人花了一块钱。但政治通道里，认同可以凭空增加或消失。一场精彩的演讲可以让一个领袖的支持率翻倍；一次丑闻可以让一个政权一夜之间失去合法性。认同不受守恒律约束——它是模因层面的东西，走的是第9章说的无压复制通道。

	权力的底座有两部分：守恒的部分（军队、税收、信息控制）和不可守恒的部分（认同、合法性、叙事）。前者是硬实力，后者是软实力。纯靠硬实力的政权是最短命的——雇佣军悖论：你付钱让人扣扳机，但钱一旦停，扳机就不扣了。枪杆子里出政权这句话对了一半：枪需要人扣，扣扳机需要认同。

	\begin{questionbox}
		「枪杆子里出政权，认同有用吗？」有用，而且比枪更根本。枪需要人扣扳机，扣扳机的人需要认同「我为什么要为这个人打仗」。历史上最持久的政权都是软硬兼施的：罗马军团（硬）+ 公民权与法律体系（软）；中华帝国（科举硬选拔）+ 天命叙事（软认同）。纯暴力的政权——从蒙古帝国到各种军事独裁——寿命都显著短于制度化的政权。
	\end{questionbox}

	\paragraph*{逻辑衔接}
	权力来自认同，但认同住在人脑里——人会死。权力怎么活得比人长？下一节回答。
	% ==================================================================
	\section{权力的模因化：制度化}
	\label{sec:pol-institutionalization}
	认同住在人脑里，人会死。一个魅力型领袖去世或失势，他的权力怎么办？答案：把个人魅力编码成程序。这就是制度化——韦伯说的从卡里斯玛到法理型的转变，用我们的语言说就是权力的模因化。

	为什么必须编码？因为人脑是易失存储，制度是不易失存储。一个将军的指挥能力住在他脑子里，他死了就没了；但一套参谋体系、一套条令条例、一套晋升规则，可以跨代传递。科层制（bureaucracy）的本质不是「官僚主义」，而是权力的冷却与压缩：把热数据（个人判断、直觉、魅力）压缩成冷数据（流程、规则、档案），牺牲灵活性换取持久性。

	宪法、礼仪、职称、军衔——这些都是编码层。它们的功能相同：让权力不依赖于任何特定个体。\lowfreq{注意}：「制度化 = 模因化」是同构类比而非恒等——模因化强调的是信息复制和传播，制度化额外加入了选择压力（哪些规则被保留、哪些被淘汰）。两者共享「从个体到集体存储」的核心机制，但制度化的选择过程更复杂。

	\begin{reader_anticipation}
		\item 编码的失真与僵化？→ 原书「结构与效率」的僵化论证在此接入：过度编码导致制度无法适应新环境，这是第12章法律通道要处理的极端情况。
	\end{reader_anticipation}

	\paragraph*{逻辑衔接}
	制度化解释了权力怎么延续，还没解释权力怎么组织。下一节把原书的集权 vs 碎片接入新位置。
	% ==================================================================
	\section{生存策略的两条扩增路径}
	\label{sec:pol-two-paths}
	原书 part2 ch2 讨论的集权 vs 碎片，在这里改写为两种不同的有压扩增组织形态。

	集权路径：降低内部交易成本。所有决策走一条线，信息集中处理，命令快速下达。优势是扩散快——政策从中央到基层只需要几层传递；劣势也是扩散快——崩溃同样沿这条线快速传播。高 $\kappa$（可预测性）意味着平时效率高，但单点故障时整个系统一起停。治水型社会天然倾向集权：洪水的应对需要协调百万人同时行动，碎片化组织做不到。

	碎片路径：冗余多、慢、抗单点故障。决策分散在多个节点，每个节点独立运作，信息通过市场信号（价格）而非行政命令协调。优势是韧性——一个节点崩了，其他节点继续跑；劣势是慢——协调成本高，响应速度低。贸易型社会天然倾向碎片：商人的网络不需要中央指挥，价格信号就够了。

	哪个更好？都不是。它们对应不同的约束条件：第4章 §4.3 的约束分析在这里变成组织形态的选择依据。高外部威胁 + 需要快速协调 → 集权；低外部威胁 + 需要创新和韧性 → 碎片。历史上成功的政治体往往是混合体：经济上碎片（市场）、军事上集权（军队）、行政上介于两者之间。

	\begin{readingnote}
		下次看到「这个国家太集权了」或「这个组织太分散了」的时候，可以问一句：它面对的主要约束是什么？如果主要约束是快速响应外部威胁，集权可能是对的；如果主要约束是适应变化的环境，碎片可能是对的。没有放之四海而皆准的组织形态。
	\end{readingnote}

	\paragraph*{逻辑衔接}
	无论哪条路径，权力都会自我扩张——这是它的扩增本性。谁来拦住它？下一章的入口。
	% ==================================================================
	\section{本章小结与衔接}
	\label{sec:pol-summary}
	政治 = 经济 + 身份权重。当交换中加入了不可货币化的偏好——人情、义务、权力——定价失效，政治通道开启。权力的来源是对职能的集体认同，认同可以凭空增减（不受守恒律约束），所以政治守恒性弱于经济。权力通过模因化（制度化）获得跨代持久性，通过集权或碎片两条路径组织自身。嵌套关系更新：文化 $\supset$ 经济 $\supset$ 政治——每条通道都是前一条加了约束的子集。

	\paragraph*{逻辑衔接}
	第12章「通道四：法律（博弈结果的编码 + 定量方法）」：当动态博弈的结果有争议、且必须给出一个数字时，法律登场。
\end{document}
